\documentclass[10pt,a4paper]{amsart}
\usepackage{amsmath,amssymb,amsthm,graphicx}
% Retained minimal preamble for the Stacks Project chapter more-morphisms.tex (More on Morphisms).
% Verbatim excerpt of the chapter preamble: Xy-pic setup (native preamble lines 14-18) and the native
% theorem declarations with the mathematical macros (native preamble lines 154-217). Copied mechanically;
% nothing added, removed or redefined. The wrapper class and the amsmath/amssymb/amsthm/graphicx base
% packages are supplied by the assembler, so the class/externaldocument part of the native preamble is omitted.
% For commutative diagrams we use Xy-pic
\usepackage[all]{xy}

% We use 2cell for 2-commutative diagrams.
\xyoption{2cell}
% Theorem environments.
%
\theoremstyle{plain}
\newtheorem{theorem}[subsection]{Theorem}
\newtheorem{proposition}[subsection]{Proposition}
\newtheorem{lemma}[subsection]{Lemma}

\theoremstyle{definition}
\newtheorem{definition}[subsection]{Definition}
\newtheorem{example}[subsection]{Example}
\newtheorem{exercise}[subsection]{Exercise}
\newtheorem{situation}[subsection]{Situation}

\theoremstyle{remark}
\newtheorem{remark}[subsection]{Remark}
\newtheorem{remarks}[subsection]{Remarks}

\numberwithin{equation}{subsection}

% Macros
%
\def\lim{\mathop{\mathrm{lim}}\nolimits}
\def\colim{\mathop{\mathrm{colim}}\nolimits}
\def\Spec{\mathop{\mathrm{Spec}}}
\def\Hom{\mathop{\mathrm{Hom}}\nolimits}
\def\Ext{\mathop{\mathrm{Ext}}\nolimits}
\def\SheafHom{\mathop{\mathcal{H}\!\mathit{om}}\nolimits}
\def\SheafExt{\mathop{\mathcal{E}\!\mathit{xt}}\nolimits}
\def\Sch{\mathit{Sch}}
\def\Mor{\mathop{\mathrm{Mor}}\nolimits}
\def\Ob{\mathop{\mathrm{Ob}}\nolimits}
\def\Sh{\mathop{\mathit{Sh}}\nolimits}
\def\NL{\mathop{N\!L}\nolimits}
\def\CH{\mathop{\mathrm{CH}}\nolimits}
\def\proetale{{pro\text{-}\acute{e}tale}}
\def\etale{{\acute{e}tale}}
\def\QCoh{\mathit{QCoh}}
\def\Ker{\mathop{\mathrm{Ker}}}
\def\Im{\mathop{\mathrm{Im}}}
\def\Coker{\mathop{\mathrm{Coker}}}
\def\Coim{\mathop{\mathrm{Coim}}}

% Boxtimes
%
\DeclareMathSymbol{\boxtimes}{\mathbin}{AMSa}{"02}

%
% Macros for moduli stacks/spaces
%
\def\QCohstack{\mathcal{QC}\!\mathit{oh}}
\def\Cohstack{\mathcal{C}\!\mathit{oh}}
\def\Spacesstack{\mathcal{S}\!\mathit{paces}}
\def\Quotfunctor{\mathrm{Quot}}
\def\Hilbfunctor{\mathrm{Hilb}}
\def\Curvesstack{\mathcal{C}\!\mathit{urves}}
\def\Polarizedstack{\mathcal{P}\!\mathit{olarized}}
\def\Complexesstack{\mathcal{C}\!\mathit{omplexes}}
% \Pic is the operator that assigns to X its picard group, usage \Pic(X)
% \Picardstack_{X/B} denotes the Picard stack of X over B
% \Picardfunctor_{X/B} denotes the Picard functor of X over B
\def\Pic{\mathop{\mathrm{Pic}}\nolimits}
\def\Picardstack{\mathcal{P}\!\mathit{ic}}
\def\Picardfunctor{\mathrm{Pic}}
\def\Deformationcategory{\mathcal{D}\!\mathit{ef}}
\newenvironment{reference}{}{}
\newenvironment{slogan}{}{}
\newenvironment{history}{}{}
\usepackage{cleveref}
\crefname{originallemma}{original lemma}{original lemmas}
\Crefname{originallemma}{Original lemma}{Original lemmas}
\makeatletter
\newlabel{algebra-lemma-quasi-finite-over-polynomial-algebra}{{00QE}{source}{}{sourcepointer-1}{}}
\newlabel{algebra-lemma-quasi-finite-over-polynomial-algebra@cref}{{[originallemma][0][]00QE}{[1][0][]source}{}{}{}}
\makeatother
\crefname{originallemma}{original lemma}{original lemmas}
\Crefname{originallemma}{Original lemma}{Original lemmas}
\makeatletter
\newlabel{algebra-lemma-where-CM}{{00RE}{source}{}{sourcepointer-2}{}}
\newlabel{algebra-lemma-where-CM@cref}{{[originallemma][0][]00RE}{[1][0][]source}{}{}{}}
\makeatother
\crefname{originallemma}{original lemma}{original lemmas}
\Crefname{originallemma}{Original lemma}{Original lemmas}
\makeatletter
\newlabel{theorem-criterion-flatness-fibre}{{theorem-criterion-flatness-fibre}{source}{}{sourcepointer-3}{}}
\newlabel{theorem-criterion-flatness-fibre@cref}{{[originallemma][0][]theorem-criterion-flatness-fibre}{[1][0][]source}{}{}{}}
\makeatother
\crefname{originallemma}{original lemma}{original lemmas}
\Crefname{originallemma}{Original lemma}{Original lemmas}
\makeatletter
\newlabel{theorem-openness-flatness}{{theorem-openness-flatness}{source}{}{sourcepointer-4}{}}
\newlabel{theorem-openness-flatness@cref}{{[originallemma][0][]theorem-openness-flatness}{[1][0][]source}{}{}{}}
\makeatother
\title{Cohen-Macaulayness at a point is equivalent to the existence of a flat quasi-finite map to affine space}
\author{The Stacks Project Authors}
\date{stacks-project:more-morphisms.tex:local-snapshot-20261011; GNU FDL 1.2 or later}
\begin{document}
\maketitle
\noindent\textit{Editorial note.} Statement span is lines 6013-6030 of sources/stacks-more-morphisms/pinned-original-source.tex; proof span is lines 6032-6065, the complete author proof (32 body lines).\par\medskip
\noindent\textit{Editorial note.} Navigation stubs: 4 cited-but-not-included labels are supplied as non-load-bearing pointers with their pinned tags where the snapshot tag index has them.\par\medskip
\noindent\textit{Editorial note.} extra\_packages is empty on purpose: the retained preamble already loads xy with options and a second declaration clashes in the assembler wrapper.\par\medskip
\noindent\textit{Editorial note.} Text is retained verbatim, including the author's own typos and placeholders.\par\medskip
\noindent\textit{Allowed standard prerequisites (editorial).}\begin{itemize}
\item Explicit prerequisite boundary. Cohen-Macaulay morphisms, quasi-finite morphisms to affine space, and the fibres' dimension; the construction is retained.
\end{itemize}\medskip
\ifdefined\hypertarget\hypertarget{sourcepointer-1}{}\fi
\noindent\textit{Source reference.} Original lemma 00QE: sources/stacks-more-morphisms/pinned-original-source.tex, source lines 6065--6065. Cited by the retained statement or proof but not part of the collected spans; supplied only as a navigation pointer. The argument does not depend on its content.\par\medskip
\ifdefined\hypertarget\hypertarget{sourcepointer-2}{}\fi
\noindent\textit{Source reference.} Original lemma 00RE: sources/stacks-more-morphisms/pinned-original-source.tex, source lines 6065--6065. Cited by the retained statement or proof but not part of the collected spans; supplied only as a navigation pointer. The argument does not depend on its content.\par\medskip
\ifdefined\hypertarget\hypertarget{sourcepointer-3}{}\fi
\noindent\textit{Source reference.} Original lemma theorem-criterion-flatness-fibre: sources/stacks-more-morphisms/pinned-original-source.tex, source lines 4590--4599. Cited by the retained statement or proof but not part of the collected spans; supplied only as a navigation pointer. The argument does not depend on its content.\par\medskip
\ifdefined\hypertarget\hypertarget{sourcepointer-4}{}\fi
\noindent\textit{Source reference.} Original lemma theorem-openness-flatness: sources/stacks-more-morphisms/pinned-original-source.tex, source lines 4446--4449. Cited by the retained statement or proof but not part of the collected spans; supplied only as a navigation pointer. The argument does not depend on its content.\par\medskip
\begin{lemma}
\label{lemma-flat-finite-presentation-characterize-CM}
Let $f : X \to S$ be a morphism of schemes which is flat and locally
of finite presentation. Let $x \in X$ with image $s \in S$.
Set $d = \dim_x(X_s)$.
The following are equivalent
\begin{enumerate}
\item $f$ is Cohen-Macaulay at $x$,
\item there exists an open neighbourhood $U \subset X$ of $x$
and a locally quasi-finite morphism $U \to \mathbf{A}^d_S$ over $S$
which is flat at $x$,
\item there exists an open neighbourhood $U \subset X$ of $x$
and a locally quasi-finite flat morphism $U \to \mathbf{A}^d_S$ over $S$,
\item for any $S$-morphism $g : U \to \mathbf{A}^d_S$
of an open neighbourhood $U \subset X$ of $x$ we have:
$g$ is quasi-finite at $x$ $\Rightarrow$ $g$ is flat at $x$.
\end{enumerate}
\end{lemma}

\begin{proof}
Openness of flatness shows (2) and (3)
are equivalent, see Theorem \ref{theorem-openness-flatness}.

\medskip\noindent
Choose affine open $U = \Spec(A) \subset X$ with $x \in U$ and
$V = \Spec(R) \subset S$ with $f(U) \subset V$. Then $R \to A$
is a flat ring map of finite presentation. Let $\mathfrak p \subset A$
be the prime ideal corresponding to $x$. After replacing $A$ by a
principal localization we may assume there exists a quasi-finite map
$R[x_1, \ldots, x_d]  \to A$, see
Algebra, Lemma \ref{algebra-lemma-quasi-finite-over-polynomial-algebra}.
Thus there exists at least one pair $(U, g)$ consisting of an
open neighbourhood $U \subset X$ of $x$ and a locally\footnote{If $S$
is quasi-separated, then $g$ will be quasi-finite.} quasi-finite morphism
$g : U \to \mathbf{A}^d_S$.

\medskip\noindent
Claim: Given $R \to A$ flat and of finite presentation, a prime
$\mathfrak p \subset A$ and $\varphi : R[x_1, \ldots, x_d] \to A$
quasi-finite at $\mathfrak p$ we have: $\Spec(\varphi)$
is flat at $\mathfrak p$ if and only if $\Spec(A) \to \Spec(R)$
is Cohen-Macaulay at $\mathfrak p$.  Namely, by
Theorem \ref{theorem-criterion-flatness-fibre}
flatness may be checked on fibres. The same is true
for being Cohen-Macaulay (as $A$ is already assumed flat over $R$).
Thus the claim follows from
Algebra, Lemma \ref{algebra-lemma-where-CM}.

\medskip\noindent
The claim shows that (1) is equivalent to (4) and combined with the
fact that we have constructed a suitable $(U, g)$ in the second
paragraph, the claim also shows that (1) is equivalent to (2).
\end{proof}

\end{document}
